% Compilation requise : XeLaTeX (latexmk -xelatex main.tex)
% Style repris du projet SCONTO-SVU (preamble.tex racine), complete par TikZ.
\usepackage[a4paper,margin=2.5cm]{geometry}
\usepackage{fontspec}
\usepackage[french]{babel}
\usepackage{microtype}
\usepackage{setspace}
\usepackage{titlesec}
\usepackage[normalem]{ulem}
\usepackage{graphicx}
\usepackage{float}
\usepackage{caption}
\usepackage{booktabs}
\usepackage{longtable}
\usepackage{tabularx}
\usepackage{array}
\usepackage{siunitx}
\usepackage{enumitem}
\usepackage{xcolor}
\usepackage{tikz}
\usetikzlibrary{arrows.meta,positioning,calc,fit,backgrounds,shapes.geometric,decorations.pathreplacing}
\usepackage{amsmath,amssymb}
\usepackage{hyperref}
\usepackage[nameinlink,noabbrev]{cleveref}
\crefname{chapter}{chapitre}{chapitres}
\Crefname{chapter}{Chapitre}{Chapitres}
\crefname{figure}{figure}{figures}
\Crefname{figure}{Figure}{Figures}
\crefname{table}{tableau}{tableaux}
\Crefname{table}{Tableau}{Tableaux}
\crefname{equation}{équation}{équations}
\Crefname{equation}{Équation}{Équations}
\usepackage{fancyhdr}
\usepackage{lastpage}
\usepackage{ifthen}

\newcolumntype{L}[1]{>{\raggedright\arraybackslash}p{#1}}
\newcolumntype{Y}{>{\raggedright\arraybackslash}X}

\IfFontExistsTF{Times New Roman}{
  \setmainfont{Times New Roman}
}{
  \setmainfont{TeX Gyre Termes}
}

\IfFontExistsTF{Arial}{
  \newfontfamily\ArialFont{Arial}
}{
  \newfontfamily\ArialFont{TeX Gyre Heros}
}

\IfFontExistsTF{Arial Black}{
  \newfontfamily\ArialBlackFont{Arial Black}
}{
  \newfontfamily\ArialBlackFont{TeX Gyre Heros}[BoldFont=* Bold]
}

% Palette sobre, utilisee uniquement pour donner un sens aux schemas.
\definecolor{ink}{RGB}{33,33,33}
\definecolor{muted}{RGB}{110,110,110}
\definecolor{paper}{RGB}{246,246,244}
\definecolor{mtsc}{RGB}{70,110,160}     % MTS / stock
\definecolor{mtoc}{RGB}{190,120,40}     % MTO / commande
\definecolor{ddmrpred}{RGB}{185,70,60}
\definecolor{ddmrpyel}{RGB}{220,180,60}
\definecolor{ddmrpgrn}{RGB}{90,150,90}
\definecolor{osc}{RGB}{110,90,150}      % Open Supply
\definecolor{resc}{RGB}{60,140,140}     % reservation

\setstretch{1.12}
\emergencystretch=1.2em
\setlength{\parindent}{0.6cm}
\setlength{\parskip}{0pt}
\setlist{nosep,leftmargin=*}

\titleformat{\chapter}[block]
  {\ArialBlackFont\bfseries\fontsize{13}{15}\selectfont}
  {\thechapter.}{0.6em}{}[\vspace{0.25ex}\titlerule]
\titlespacing*{\chapter}{0pt}{-10pt}{16pt}

\titleformat{\section}
  {\ArialFont\bfseries\fontsize{11}{13}\selectfont}
  {\thesection}{0.6em}{}
\titleformat{\subsection}
  {\ArialFont\bfseries\fontsize{11}{13}\selectfont}
  {\thesubsection}{0.6em}{}

\captionsetup{font=small,labelfont=bf,justification=centering}
\sisetup{locale=FR,detect-all}

\hypersetup{
  colorlinks=true,
  linkcolor=black,
  citecolor=black,
  urlcolor=black,
  pdftitle={Intégration des stratégies MTO et DDMRP dans SCONTO-SVU},
  pdfauthor={}
}

\pagestyle{fancy}
\fancyhf{}
\fancyhead[L]{\small SCONTO-SVU : MTO et DDMRP}
\fancyhead[R]{\small Présentation technique et validation}
\fancyfoot[C]{\small Page \thepage\ sur \pageref{LastPage}}
\renewcommand{\headrulewidth}{0.3pt}
\renewcommand{\footrulewidth}{0pt}

% Terme technique : code Java en police machine, premiere occurrence expliquee dans le texte.
\newcommand{\jv}[1]{\texttt{#1}}
\newcommand{\Etat}[1]{\texttt{#1}}
